\documentclass[12pt]{amsproc}
\usepackage{mynotes}

\renewcommand{\course}{Hopf algebras}

% special macros
\newcommand{\ydH}{\prescript{H}{H}{\mathcal{YD}}}
\newcommand{\lmod}[1]{\prescript{}{#1}{\mathcal{M}}}
\newcommand{\rmod}[1]{\mathcal{M}_{#1}}
\newcommand{\lcomod}[1]{\prescript{#1}{}{\mathcal{M}}}
\newcommand{\rcomod}[1]{\mathcal{M}^{#1}}



\title{\course}
\author{Leandro Vendramin}
\address{Department of Mathematics and Data
Science, Vrije Universiteit Brussel, Pleinlaan 2, 1050 Brussel}
\email{Leandro.Vendramin@vub.be}
\thanks{}
\date{}

\begin{document}

\begin{abstract}
    The notes correspond to the bachelor 
    course \textbf{Hopf algebras} of the 
    Vrije Universiteit Brussel, 
    Faculty of Sciences, 
    Department of Mathematics and Data Sciences. 
\end{abstract}

\maketitle

%\begin{figure}[h]
%    \includegraphics[scale=0.5]{hungerford.png}
%\end{figure}

\thispagestyle{empty}

\setcounter{tocdepth}{1}

\newpage
\thispagestyle{empty}
\tableofcontents

\newpage
\pagenumbering{arabic}

\include{introduction}

\input{01}
\input{02}
\input{03}
\input{04}

\input{solutions}
\bibliographystyle{abbrv}
\bibliography{refs}

\printindex

\end{document}
